%!TEX root = ../main.tex

\textbf{Актуальність теми.} Цифрові музичні колекції, потокові сервіси та
архіви звукозаписів містять мільйони творів, для більшості з яких відсутній
опис інструментального складу. Інформація про те, які інструменти звучать у
записі, потрібна для каталогізації та пошуку музики, побудови рекомендаційних
систем, музикознавчого аналізу й навчальних застосунків. Ручне анотування
великих колекцій потребує значних витрат часу фахівців, тому актуальною є
задача автоматичного розпізнавання музичних інструментів у звукових записах.

Найскладнішим є випадок багатоінструментних записів, у яких кілька
інструментів звучать одночасно, їхні спектри накладаються, а набір
присутніх інструментів заздалегідь невідомий. Сучасні підходи на основі
згорткових нейронних мереж суттєво перевершили класичні методи, проте
більшість із них використовують архітектури, запозичені з обробки
зображень, і навчаються на записах з одним переважаючим інструментом, тоді
як застосовуються до багатоінструментних. Крім того, методики оцінювання в
опублікованих роботах часто не описані повністю або передбачають
підбирання налаштувань на тестових даних, що ускладнює чесне порівняння
методів. Отже, актуальним є розроблення моделі, яка враховує тембральну та
часову будову звуку, навчається з урахуванням багатоінструментного
характеру музики та оцінюється за відтворюваною методикою.

\textbf{Метою дослідження} є підвищення точності розпізнавання музичних
інструментів у багатоінструментних звукових записах шляхом розроблення
моделі згорткової нейронної мережі з музично обґрунтованими фільтрами та
методу її навчання на синтетичних сумішах записів.

Для досягнення мети необхідно розв'язати такі \textit{задачі}:
\begin{enumerate}
    \item проаналізувати існуючі методи автоматичного розпізнавання музичних
    інструментів, набори даних і методики оцінювання;
    \item формалізувати задачу багатоміткового розпізнавання інструментів та
    розробити математичну модель розпізнавання: часово-частотне представлення,
    музично обґрунтовані згорткові фільтри, синтез багатоінструментних
    сумішей, правило прийняття рішення;
    \item розробити програмне забезпечення для підготовки даних, навчання,
    розпізнавання та оцінювання моделей;
    \item розробити відтворювану методику оцінювання та провести валідацію
    розробленої моделі на стандартизованому наборі даних, порівнявши її з
    існуючими рішеннями та дослідивши внесок кожної складової моделі.
\end{enumerate}

\textbf{Об'єктом дослідження} є процес автоматичного розпізнавання
музичних інструментів у багатоінструментних звукових записах на основі
згорткових нейронних мереж.

\textbf{Предметом дослідження} є моделі та методи побудови згорткових
нейронних мереж для розпізнавання музичних інструментів, що враховують
тембральну та часову будову звуку.

\textbf{Методи дослідження}: методи цифрової обробки звукових сигналів
(коротковіконне перетворення Фур'є, мел-спектральний аналіз), методи
навчання згорткових нейронних мереж, синтез багатоінструментних сумішей для
розширення навчальних даних та порівняльний експериментальний аналіз
моделей на стандартизованому наборі даних IRMAS з оцінюванням якості за
макро- та мікроусередненою F1-мірою.

\textbf{Наукова новизна} одержаних результатів полягає в такому:
\begin{enumerate}
    \item показано, що на мел-шкалі одна октава займає від трьох до
    двадцяти восьми частотних смуг залежно від висоти, через що вертикальні
    (тембральні) згорткові фільтри сталої висоти охоплюють різні музичні
    інтервали на різних частотах; запропоновано виконувати такі фільтри на
    логарифмічній частотній осі, де кожен фільтр охоплює однаковий музичний
    інтервал незалежно від висоти ноти;
    \item встановлено, що користь явного гармонічного представлення
    (накладання каналів на частотах обертонів) залежить від умов навчання:
    воно покращує розпізнавання при навчанні на записах з одним
    інструментом, але погіршує його при навчанні на синтетичних
    багатоінструментних сумішах;
    \item розроблено модель, яка поєднує музично обґрунтовані фільтри з
    навчанням на синтетичних сумішах і за суворішої методики оцінювання
    перевершує опубліковані результати відомих згорткових моделей на наборі
    даних IRMAS за меншої кількості параметрів.
\end{enumerate}

\textbf{Практичне значення} одержаних результатів полягає в розробленому
програмному забезпеченні, яке може бути використане для автоматичної
каталогізації та пошуку в музичних колекціях, а також як основа для
подальших досліджень тембрового аналізу та змістового музичного пошуку.